\paragraph{Hint 1.} Quantify over an arbitrary test scheme $T$.
\paragraph{Hint 2.} Universal properties are invariant under compatible isomorphism.
\paragraph{Hint 3.} The second component is forced to be the structure map of the first.
\paragraph{Hint 4.} Swap projections.
\paragraph{Hint 5.} Tensor the polynomial rings over $k$.
\paragraph{Hint 6.} Use $(A/I)\otimes_A(A/J)\cong A/(I+J)$.
\paragraph{Hint 7.} Use $B\otimes_AA_f\cong B_{f}$.
\paragraph{Hint 8.} There is exactly one map from the empty scheme to any scheme.
\paragraph{Hint 9.} Use multiplication $B\otimes_AB\to B$.
\paragraph{Hint 10.} Send $b\otimes c$ to $b f^\#(c)$.
\paragraph{Hint 11.} This is part of the defining commutative square.
\paragraph{Hint 12.} The twice-swapped map has the original projections.
\paragraph{Hint 13.} Compare the functors represented by both bracketings.
\paragraph{Hint 14.} Combine all polynomial variables.
\paragraph{Hint 15.} Add the three ideals.
\paragraph{Hint 16.} The diagonal represents equality.
\paragraph{Hint 17.} The graph represents the equation $y=f(x)$.
\paragraph{Hint 18.} Look at primes of a residue-field tensor product.
\paragraph{Hint 19.} Use a separable field extension that splits after scalar extension.
\paragraph{Hint 20.} Use the tensor-product universal property.
\paragraph{Hint 21.} Localization is a tensor product and tensor products associate.
\paragraph{Hint 22.} Tensor the two polynomial rings over $\mathbb Z$.
\paragraph{Hint 23.} Two overlap identifications with the same projections must coincide.
\paragraph{Hint 24.} Write the natural Hom bijection.
\paragraph{Hint 25.} Only identify the pullback here; defer geometric analysis.
\paragraph{Hint 26.} Their migration destination is explicitly VI/24 in the ledger.